\chapter*{Résumé}
\phantomsection
\addcontentsline{toc}{chapter}{Résumé}

Ce projet de fin d'études a été réalisé au sein de SOMASTEEL, une entreprise marocaine spécialisée dans la sidérurgie, le laminage à chaud des aciers et la production de fer à béton. Le projet vise à répondre à un besoin du service logistique concernant le suivi des camions loués transportant la ferraille depuis le port vers l'usine.

Contrairement aux camions appartenant à l'entreprise, les camions loués ne peuvent pas être équipés de dispositifs GPS. Avant la mise en place de la solution, leur suivi reposait principalement sur des échanges téléphoniques, ce qui limitait la précision, la centralisation et la fiabilité des informations.

La solution proposée, intitulée \textbf{SomaScan}, repose sur l'identification des camions par des codes QR associés aux plaques d'immatriculation. Elle permet aux opérateurs du port et de l'usine de scanner les camions à travers une application mobile développée avec Flutter, tandis qu'un agent logistique assure la gestion et le suivi des expéditions à travers une application web développée avec React. Le backend, développé avec Laravel, expose des API REST pour gérer les données, l'authentification, les rôles et les statuts des camions.

Le projet couvre l'analyse du besoin, la conception, le développement, le déploiement ainsi que la proposition d'un workflow DevOps destiné à automatiser et fiabiliser les futures mises en production.

\vspace{0.5cm}

\noindent\textbf{Mots-clés :} SOMASTEEL, SomaScan, logistique, traçabilité, camions loués, ferraille, QR Code, Laravel, Flutter, React, API REST, DevOps.
